%% C4 nivel 3 — Diagrama de componentes (plantilla)
%% Compilar: tectonic -X compile c4-componentes.tex
\documentclass[11pt,a4paper]{article}
\usepackage{../../docs/latex/legasoft}
\usepackage{../../docs/latex/legasoft-diagramas}

\lgtitulo{C4 nivel 3: diagrama de componentes}
\lgtituloCorto{C4 — Componentes}
\lgcodigo{LGS-MOD-C4-3}
\lgversion{0.1}
\lgestado{Plantilla}
\lgfecha{\campo{fecha}}
\lgautor{\lgArquitecto\ (\lgArquitectoRol)}

\begin{document}

\begin{center}
{\sffamily\bfseries\LARGE\color{lgPrimario} C4 nivel 3: diagrama de componentes\par}
\vspace{2pt}
{\sffamily\normalsize\color{lgGris} Contenedor \campo{API de aplicación} · Legasoft\par}
\end{center}
\vspace{6pt}
{\color{lgBorde}\hrule height 0.8pt}
\vspace{10pt}

\begin{porcompletar}[Cómo usar esta plantilla]
El nivel 3 abre \emph{un} contenedor y muestra sus componentes internos y las
responsabilidades de cada uno. Se elabora solo para los contenedores cuya
estructura interna necesite explicarse; no es obligatorio hacerlo para todos.
\end{porcompletar}

% =============================================================================
\section*{Diagrama}

\begin{figure}[H]
\centering
\ajustaancho{%
\begin{tikzpicture}[c4]

  \node[c4contenedor, text width=3.4cm, align=center, minimum height=1.5cm] (web)
    {\ctexto{Aplicación web}{\campo{tecnología}}{Consume la API.}};

  % --- Componentes internos -------------------------------------------------
  \node[c4componente, below=2.0cm of web] (controlador)
    {\ctexto{Controladores}{}{Reciben las peticiones y validan la entrada.}};

  \node[c4componente, below=1.5cm of controlador] (servicio)
    {\ctexto{Servicios de dominio}{}{Concentran las reglas de negocio.}};

  \node[c4componente, below=1.5cm of servicio] (repo)
    {\ctexto{Repositorios}{}{Aíslan el acceso a la persistencia.}};

  \node[c4componente, right=2.2cm of servicio] (seguridad)
    {\ctexto{Seguridad}{}{Autenticación, autorización y auditoría.}};

  % Límite del contenedor
  \begin{scope}[on background layer]
    \node[c4limite, fit=(controlador)(servicio)(repo)(seguridad),
          label={[font=\sffamily\scriptsize\itshape, text=lgPrimario,
                  anchor=north west]north west:{Contenedor: \campo{API de aplicación}}}]
      {};
  \end{scope}

  \node[c4datos, below=1.6cm of repo, text width=2.6cm, align=center,
        minimum height=1.7cm] (bd)
    {\ctexto{Base de datos}{\campo{gestor}}{Persistencia.}};

  % --- Relaciones -----------------------------------------------------------
  \draw[c4rel] (web.south) -- (controlador.north)
    node[c4etq, pos=0.5] {\scriptsize[JSON/HTTPS]};
  \draw[c4rel] (controlador.south) -- (servicio.north) node[c4etq, pos=0.5] {invoca};
  \draw[c4rel] (servicio.south) -- (repo.north) node[c4etq, pos=0.5] {consulta};
  \draw[c4rel] (repo.south) -- (bd.north) node[c4etq, pos=0.5] {\scriptsize\campo{protocolo}};
  \draw[c4rel] (controlador.east) to[out=0, in=90] (seguridad.north);
  \draw[c4rel] (servicio.east) -- (seguridad.west) node[c4etq, pos=0.5] {verifica};

\end{tikzpicture}}
\caption{Componentes internos del contenedor y sus dependencias.}
\end{figure}

% =============================================================================
\Needspace*{20\baselineskip}
\section*{Componentes}

\begin{center}
\begin{tabularx}{\linewidth}{@{}L{3.6cm} Y@{}}
\toprule
\sffamily\bfseries\color{lgPrimario}Componente &
\sffamily\bfseries\color{lgPrimario}Responsabilidad \\
\midrule
Controladores        & Traducen peticiones externas al modelo interno y validan la entrada. No contienen reglas de negocio. \\
Servicios de dominio & Concentran las reglas de negocio y coordinan las operaciones. \\
Repositorios         & Aíslan la persistencia para que las reglas de negocio no dependan del gestor de datos. \\
Seguridad            & Autenticación, autorización y registro de auditoría. \\
\bottomrule
\end{tabularx}
\captionof{table}{Responsabilidad de cada componente.}
\end{center}

\begin{nota}[Criterio de diseño]
La separación en capas responde al atributo de mantenibilidad. Si el proyecto no
justifica esta estructura, simplifíquela: Legasoft evita introducir complejidad
sin un beneficio verificable.
\end{nota}

\end{document}
